%---------------------------------------------------------------------------%
%->> 封面信息及生成
%---------------------------------------------------------------------------%
%-
%-> 中文封面信息
%-
\confidential{}% 密级：只有涉密论文才填写
\schoollogo{scale=0.2}{ShanghaiTechLogo}% 校徽
%======论文题目 页眉显示 务必准确========
\title{灵巧操作的数据采集与策略部署}% 论文中文题目
\author{黄晟真}% 论文作者
\ID{2022531025}
\entranceYear{2022}
\advisor{陈嘉豪}% 指导教师：姓名
\advisorsec{}% 第二指导老师：按情况填写
\degree{学士}% 学位：学士、硕士、博士
\degreetype{工学}% 学位类别：理学、工学、工程、医学等
\major{电子信息工程}% 二级学科专业名称
\institute{信息科学与技术学院}% 院系名称
\chinesedate{二零二六年六月}% 毕业日期：夏季为6月、冬季为12月
%-
%-> 英文封面信息
%-
\englishtitle{Data Collection and Strategy Deployment for Dexterous Manipulation}% 论文英文题目
\englishauthor{Huang Shengzhen}% 论文作者
\englishadvisor{Chen Jiahao}% 指导教师
\englishdegree{Bachelor}% 学位：Bachelor, Master, Doctor。封面格式将根据英文学位名称自动切换，请确保拼写准确无误
\englishdegreetype{Engineering}% 学位类别：Philosophy, Natural Science, Engineering, Economics, Agriculture 等
\englishthesistype{thesis}% 论文类型： thesis, dissertation
\englishmajor{Electronic Engineering}% 二级学科专业名称
\englishinstitute{School of Information Science and Technology}% 院系名称
\englishdate{June 2026}% 毕业日期：夏季为June、冬季为December
%-
%-> 生成封面
%-
%-======================================================
% 封面和声明页格式不符合教务处要求，因此此处不生成pdf格式的 ||
% 论文封面和声明页，请各位使用word模板中的相关内容。       ||
%=======================================================
%\maketitle% 生成中文封面
%\makeenglishtitle% 生成英文封面
%-
%-> 作者声明
%-
%\makedeclaration% 生成声明页
%-
%-> 中文摘要
%-

\chapter*{摘\quad 要}\chaptermark{摘\quad 要}% 摘要标题
\setcounter{page}{1}% 开始页码
\pagenumbering{Roman}% 页码符号

灵巧操作是具身智能中最具挑战性的研究方向之一，其关键难点在于高自由度本体控制、复杂接触交互以及真实示教数据的高成本获取。针对上述问题，本文设计并实现了一套面向灵巧操作任务的数据采集与策略部署流程，构建了由天机 Marvin 机械臂、WUJI HAND 灵巧手、HTC Vive Tracker 动捕系统、MANUS 数据手套和多视角相机组成的遥操作实验平台，并基于 ROS 2 完成多模态数据同步、记录与管理。

在数据采集与重定向方面，本文建立了人体动作到机器人动作的映射方法，将操作者的上肢位姿和手部骨架分别转换为机械臂关节角与灵巧手关节位置，并形成可用于策略学习的专家示教数据集。在此基础上，本文进一步以 $\pi_{0.5}$ 基座模型为核心，对灵巧操作任务进行真机微调与部署。实验中，模型以关节空间作为动作表示，通过对右臂关节增量和灵巧手绝对关节位置的联合建模，实现了从视觉观测和语言指令到连续动作块的端到端预测。

在真机验证中，微调后的策略已经能够在部分试次中完成“将紫色圆柱体放入碗内”的任务，但整体成功率仍然有限。失败现象主要表现为抓取位置错误、手部过低引发桌面碰撞以及动作输出退化。分析表明，造成性能受限的原因可能包括示教数据中阻抗控制与负载补偿带来的姿态偏差，以及实验环境光照不足导致的视觉识别不稳定。本文的工作验证了高质量遥操作数据对灵巧操作策略学习的重要性，也为后续进一步提升真机部署鲁棒性提供了基础。

\keywords{灵巧操作，数据采集，遥操作，动作重定向，视觉-语言-动作模型}% 中文关键词
%-
%-> 英文摘要
%-
\chapter*{Abstract}\chaptermark{Abstract}% 摘要标题

Dexterous manipulation remains one of the most challenging problems in embodied intelligence, mainly due to the high-dimensional control space, complex contact-rich interactions, and the high cost of collecting high-quality expert demonstrations. To address these issues, this thesis develops an integrated pipeline for dexterous manipulation data collection and policy deployment. The proposed platform combines a Tianji Marvin arm, a WUJI HAND dexterous hand, an HTC Vive Tracker motion-capture system, MANUS data gloves, and multi-view RGB cameras, and uses ROS 2 for multimodal synchronization, logging, and session management.

For data collection and retargeting, the thesis establishes a mapping from human motion to robot actions, converting the operator's upper-limb pose and hand skeleton into robot arm joint angles and hand joint positions. The resulting demonstrations are then used to fine-tune and deploy a $\pi_{0.5}$ base model on the real robot. In the deployed policy, actions are represented in joint space, with the 7 arm joints modeled as relative deltas and the 20 hand joints kept as absolute positions, enabling end-to-end prediction from visual observations and language instructions to continuous action chunks.

In real-robot experiments, the fine-tuned policy can complete the task of placing a purple cylinder into a bowl in part of the trials, but the overall success rate remains limited. The main failure modes include incorrect grasp locations, the hand moving too low and colliding with the table, and degenerated actions with little or no motion. Analysis suggests that the performance bottleneck is likely caused by pose bias in the demonstration data due to impedance control and incomplete end-effector compensation, as well as unstable visual recognition under poor lighting conditions. Overall, this work demonstrates the importance of high-quality teleoperation data for dexterous manipulation policy learning and provides a foundation for more robust real-robot deployment in future work.

\englishkeywords{dexterous manipulation, data collection, teleoperation, action retargeting, vision-language-action model}% 英文关键词
%---------------------------------------------------------------------------%
